\begin{lemma}[A terminal affine-$B$ tail]
\label{lem:terminal-b-tail}
Let $K$ be a nondegenerate Lorentzian subspace of the physical root space.
Suppose that a root $e$ of squared norm $2$ has orthogonal residual
$e_\perp$ with $e_\perp^2=\rho<1$. There cannot be a chain
$f_1,\ldots,f_k$ of roots in $K^\perp$, each of squared norm $2$,
with the following native products. For $k\ge2$, consecutive products
are $-1$, except for $(f_{k-1},f_k)=-\sqrt2$; also
$(e,f_1)=-1$. For $k=1$, instead $(e,f_1)=-\sqrt2$.
All other products between distinct roots in the chain, and between
$e$ and $f_j$ for $j>1$, are zero.
\end{lemma}

\begin{proof}
The complement $K^\perp$ is positive definite. Put
\[
 v=f_1+\cdots+f_{k-1}+\frac{f_k}{\sqrt2}.
\]
For $k\ge2$, its diagonal terms sum to $2(k-1)+1$ and its
consecutive cross terms sum to $-2(k-1)$. Thus $v^2=1$.
For $k=1$ the same identity follows from $f_1^2=2$.
The prescribed products give $(e,v)=-1$ in both cases.
Since $v\in K^\perp$, positive-space Cauchy gives the contradiction
\[
 1=(e_\perp,v)^2\le e_\perp^2v^2=\rho<1.
\]
Only the coefficient of the final root has changed: every physical
root still has squared norm $2$.
\end{proof}

This lemma treats a terminal chain of any length, including cases where
the first two new roots alone do not violate Cauchy. Its orthogonality
hypothesis must be proved before it is used. In our applications the
complete contact certificate permits a unique first root, or permits
that root together with one exterior root. In the latter case the
full no-bridge estimate bounds the whole component in the presence of
the exterior root; that bound is strictly smaller than the actual target
order. This excludes the exterior root and every indirect bridge.
First-contact uniqueness then proves successively that all later roots
are genuinely new and orthogonal to the old core. Their native zero
products can be used only after this step. If the forced first contact
is at $i$ with magnitude $1$, the inverse certificate gives
$\rho=2-J_{ii}<1$.

The exact application checks the complete supplied terminal-suffix map,
all its edges and nonedges, and the original physical contact-profile
selector. A contact theorem computed for another target is used only
when its original core, ordered corner and full contact hypotheses
coincide; its old target map is never relabelled as the present one.
The terminal congruence in the calculation of $v^2$ is a rational check,
with the scale $1/\sqrt2$ recorded explicitly.

